\section*{Planned Demonstration}

The finished project will be a two-dimensional particle simulation written in C++. This particle
simulation will show how spatial partitioning can make collision-detection faster. Particles will move
inside a space collide with each other and bounce off the walls. Users can change how particles there
are change their spatial distribution, pause or resume the particle simulation and choose between three
collision-detection algorithms: brute force, a uniform spatial grid and a Quadtree.

The application will show a real-time visualization of the Quadtree's subdivisions. This
visualization will let users see how the environment is split into regions as particle density grows. When
a collision query runs the application will highlight the Quadtree regions and the candidate particle
pairs. This makes it clear how the Quadtree limits the search to particles instead of checking every pair.
A performance dashboard will show metrics like candidate-pair counts, collision-detection time,
index-construction time and the number of Quadtree nodes.

During the demonstration I will start with a small number of particles and show how each
collision-detection algorithm finds collisions. Then I will gradually increase the particle count to show the
cost of brute-force collision detection compared to the spatial partitioning methods. Next I will present both
uniform and clustered particle distributions to illustrate how workload patterns affect the
Quadtree's performance versus the uniform grid. Finally, I will run the application's benchmarking mode
with particle setups and display the resulting performance measurements.

The main goal of the demonstration is to make the Quadtree's advantage clear both visually and
in numbers. The spatial subdivision overlay will show how the Quadtree organizes the environment. At
the time the performance dashboard will reveal whether this organization cuts down on needless
collision checks and speeds up execution. The demonstration will also point out cases where building
and walking through the Quadtree adds work or where dense clusters reduce its usefulness. By mixing
visualization, with measurable performance data the finished project will show not only how the
Quadtree works but also when and why it can make collision detection more efficient.
